\documentclass[11pt, a4paper]{article}

\usepackage[margin=2.5cm]{geometry}
\usepackage[T1]{fontenc}
\usepackage[utf8]{inputenc}
\usepackage{microtype}
\usepackage{parskip}
\usepackage{titlesec}
\usepackage{enumitem}
\usepackage{booktabs}
\usepackage{longtable}
\usepackage{array}
\usepackage{xcolor}
\usepackage{listings}
\usepackage{fancyhdr}
\usepackage[hidelinks]{hyperref}
\usepackage{amsmath}
\usepackage{graphicx}

% header / footer
\pagestyle{fancy}
\fancyhf{}
\fancyhead[L]{PlayForge Manager}
\fancyhead[R]{Milestone \#2}
\fancyfoot[C]{\thepage}
\renewcommand{\headrulewidth}{0.4pt}

% inline code
\definecolor{codegray}{gray}{0.95}
\newcommand{\code}[1]{\texttt{#1}}

% listings styling for Java snippets
\lstset{
    basicstyle=\ttfamily\footnotesize,
    backgroundcolor=\color{codegray},
    frame=single,
    breaklines=true,
    columns=fullflexible,
    keepspaces=true,
    showstringspaces=false,
    tabsize=4
}

\title{}
\author{}
\date{}

\begin{document}

% ========================= TITLE PAGE =========================
\begin{titlepage}
    \centering
    \vspace*{2cm}

    {\large CE 216 -- Fundamental Topics in Programming \par}
    \vspace{0.4cm}
    {\normalsize Associate Prof.\ Dr.\ Kaya Oğuz \par}
    \vspace{0.8cm}

    \IfFileExists{logo.png}{%
        \includegraphics[width=0.44\textwidth]{logo.png}\par
        \vspace{0.8cm}
    }{}

    \vfill

    {\Huge \textbf{PlayForge Manager} \par}
    \vspace{0.6cm}
    {\Large Milestone \#2 Report \par}
    \vspace{0.4cm}
    {\large Implementation, Testing, and Design Changes \par}

    \vfill

    \begin{tabular}{ll}
        \textbf{Team:}          & JavaIsSlow                                                           \\
        \textbf{Members:}       & Mirza İşi, Sepehr Rahimi, Doruk Faydalı, Çağan Parlapan              \\
        \textbf{Repository:}    & \url{https://github.com/mirzaisi/CE216_Project_Team_JavaIsSlow}      \\
        \textbf{Submission:}    & 12.04.2026                                                           \\
    \end{tabular}

    \vfill
    {\large Spring 2026 \par}
\end{titlepage}

% ========================= TOC =========================
\tableofcontents
\newpage

% ========================= 1. SCOPE =========================
\section{Milestone Scope and Deliverables}

This document accompanies the Milestone \#2 submission of PlayForge Manager.
The goal of M2 is to translate the M1 design into a working Java project
with a complete abstraction layer, a first concrete sport,
a minimal runnable simulation, and a focused unit-test suite that proves
the foundations are sound.

The expected deliverables, following the assignment instructions, are:
\begin{itemize}
    \item A Maven project that builds and tests cleanly from a fresh clone.
    \item A runnable entry point reachable via \code{mvn exec:java}.
    \item A minimum of twenty meaningful unit tests. The project ships with
          \textbf{88}, covering every major abstraction and policy.
    \item This milestone report (PDF) describing what was built and what
          changed from the M1 design.
    \item A GitHub Release created from the final submission commit.
    \item A TXT file containing the repository URL.
\end{itemize}

% ========================= 2. SUMMARY =========================
\section{Summary of Implementation}

Every interface and abstract class defined in the M1 design document is
implemented in code. Football is realised as a first concrete sport that
plugs into the shared abstractions without leaking back into the core layer.

\subsection{Layered Package Layout}

The project follows the five-layer structure defined in Section 4 of the
M1 design:

\begin{itemize}
    \item \code{com.playforgemanager.core} -- shared abstractions, interfaces,
          and abstract base classes. No football imports.
    \item \code{com.playforgemanager.application} -- use-case services
          (initialization, session wiring). Depends only on core.
    \item \code{com.playforgemanager.football} -- concrete football sport
          module implementing every core contract.
    \item \code{com.playforgemanager.infrastructure} -- support services,
          currently an in-memory asset provider for names, teams, and logos.
    \item \code{com.playforgemanager.main} -- program entry point used by
          \code{mvn exec:java}.
\end{itemize}

The dependency rule is enforced by the package layout and verified by a
simple import check: neither \code{core} nor \code{application} imports
anything from the football module.

\subsection{Shared Abstractions (Core Layer)}

All nineteen abstractions listed in Table~5.1 of the M1 design are
implemented. The code also defines \code{ProgressionState} as a small
implementation enum used by \code{GameSession}; it is not counted as one
of the Table~5.1 abstractions.

\begin{longtable}{ll}
    \toprule
    \textbf{Abstraction}      & \textbf{Form}                          \\
    \midrule
    \endhead
    \code{GameSession}        & concrete class                         \\
    \code{Sport}              & interface                              \\
    \code{SportFactory}       & interface                              \\
    \code{Ruleset}            & interface                              \\
    \code{League}             & abstract class                         \\
    \code{Season}             & abstract class                         \\
    \code{Team}               & abstract class                         \\
    \code{Player}             & abstract class                         \\
    \code{Coach}              & abstract class                         \\
    \code{Match}              & abstract class                         \\
    \code{Fixture}            & concrete class                         \\
    \code{Lineup}             & interface                              \\
    \code{Tactic}             & interface                              \\
    \code{TrainingPlan}       & interface                              \\
    \code{MatchEngine}        & interface                              \\
    \code{Scheduler}          & interface                              \\
    \code{StandingsPolicy}    & interface                              \\
    \code{InjuryPolicy}       & interface                              \\
    \code{AssetProvider}      & interface                              \\
    \bottomrule
\end{longtable}

\subsection{Football Module}

The football package provides concrete implementations for the football
roles from the M1 design. The scheduler role is implemented by
\code{RoundRobinFootballScheduler}, a more explicit name for the double
round-robin scheduler described in the design. The module contains:

\begin{itemize}
    \item \code{FootballSport}, \code{FootballSportFactory} -- composition
          root and creation entry point.
    \item \code{FootballRuleset} -- scoring, the fixed 1 goalkeeper /
          4 defender / 4 midfielder / 2 forward starting-lineup policy,
          bench size, points logic, and availability checks.
    \item \code{RoundRobinFootballScheduler} -- double round-robin fixture
          generator.
    \item \code{FootballLeague}, \code{FootballSeason} -- competition
          containers and week-by-week progression.
    \item \code{FootballStandingsPolicy}, \code{FootballStandingRow} --
          recorded-match bookkeeping and league table calculation with
          tie-breaking.
    \item \code{FootballMatch}, \code{FootballMatchEngine} -- match entity
          and the stateless engine that produces results.
    \item \code{FootballInjuryPolicy} -- post-match injury application and
          future-match recovery countdown.
    \item \code{FootballTeam}, \code{FootballPlayer}, \code{FootballCoach},
          \code{FootballLineup}, 
          \code{FootballTactic}, \code{FootballTrainingPlan}--football realisations of the core entities.
    \item \code{FootballAttributeProfile}, \code{FootballPosition} --
          composed pieces of football-specific player data.
    \item \code{FootballConsoleDemo} -- thin console runner used by
          \code{Main}.
\end{itemize}

\subsection{Application and Infrastructure}

\code{GameInitializationService} is the only application service required
for M2. It accepts any \code{SportFactory}, creates a sport, builds a
league, generates a fixture list, and returns a fully wired
\code{GameSession}. The service depends only on the core layer.

\code{InMemoryAssetProvider} supplies male and female given names, team
names, and placeholder logo references. It is used by the football factory
during league and team generation.

% ========================= 3. CHANGELOG =========================
\section{Changes from the M1 Design}

The M1 design was close to what we actually built, but a few adjustments
were made during implementation. They are listed here with the reason for
each change.

\subsection{Bootstrap layer removed}

Early in M2 we kept a parallel \code{Bootstrap*} family of football
classes alongside the canonical ones. These had been created during the
M1 transition so the project always compiled, and were used by some
tests as lightweight stand-ins for the real types.

But later, these nine classes were deleted
(\code{BootstrapFootballLeague}, \code{Season}, \code{Team}, \code{Player},
\code{Coach}, \code{Lineup}, \code{Match}, \code{Scheduler},
\code{Tactic}), the console demo was renamed from
\code{BootstrapFootballConsoleDemo} to \code{FootballConsoleDemo}, and the
tests that referenced the bootstrap types were migrated onto the canonical
football classes. The reason was to remove duplication and make the code
match the M1 football roles directly. The scheduler implementation keeps
the more specific name \code{RoundRobinFootballScheduler} because M2 fixes
the league format to a double round-robin.

\subsection{FootballMatchEngine and FootballInjuryPolicy extracted}

The M1 design lists \code{FootballMatchEngine} and \code{FootballInjuryPolicy}
as top-level classes. During the initial M2 prototyping they were inlined as
private static inner classes inside \code{FootballSport}, which worked but
obscured the design. Both were extracted into their own files so that
\code{FootballSport} is now a thin composition root and each policy can be
reused, tested, and replaced independently.

\subsection{Deterministic match simulation}

The M1 design describes the match engine as stateless but does not
prescribe how it should produce scores. For M2 we chose a simple
deterministic simulation: the engine derives a seed from the participating
team names and feeds it to a local \code{java.util.Random}. This choice is
intentional -- it gives us a reproducible result for each team pairing,
which makes the simulation trivial to unit-test and makes the console demo
stable across runs. Player attributes and tactics are stored on the entities
but are not yet folded into the scoring formula. Integrating them is on the
list of deferred work for later milestones.

\subsection{Injury model simplified}

The M1 design lists \code{InjuryPolicy} as a policy that handles injuries
and recovery but does not specify the exact rule set. For M2 we chose a
minimal, deterministic model:

\begin{itemize}
    \item After every played match the losing side has one player from the
          starting lineup marked as injured. On a draw the home side takes
          the hit. The player is selected by a deterministic index derived
          from the final scoreline so the rule is reproducible in tests.
    \item Injured players are unavailable for a configurable number of
          future match opportunities (default: two). A player injured in
          the current played week is not decremented by that same week's
          recovery pass; the countdown ticks down after later missed
          match weeks.
    \item The football ruleset rejects any lineup that contains an
          unavailable player, so injured players cannot be selected.
    \item When a player is injured, the affected team's stored lineup is
          cleared so that the next weekly resolution rebuilds a fresh
          valid lineup from whoever is still available. This is the
          explicit link between injury and availability that connects the
          policy to the season flow.
\end{itemize}

This model is documented here as a deliberate simplification.

\subsection{Football lineup policy made explicit}

This code uses one deterministic starting-lineup composition everywhere:
1 goalkeeper, 4 defenders, 4 midfielders, and 2 forwards. The football
factory and the season's automatic lineup resolution both build this
shape through \code{FootballRuleset}, and the ruleset rejects starting
elevens with extra goalkeepers, missing forwards, or any other positional
composition. Generated squads contain nineteen players so a two-match
injury window has enough reserve depth by position.

\subsection{Standings and core invariants tightened}

\code{FootballStandingsPolicy.recordMatch} now records played matches so
standings can change through the policy method as well as through fixture
recomputation. The table recomputation path skips duplicate match objects,
so a scheduled match recorded through both paths is counted once. The core
\code{League} and \code{Fixture} classes now reject duplicate team IDs,
fixtures with non-member teams, and attempts to attach an unplayed match.

\subsection{Clean-up of FootballSeason.resolveLineup}

The initial version of \code{FootballSeason.resolveLineup} contained
dead branches: after narrowing the team to a \code{FootballTeam}, a second
generic branch built a list of available players that was never used and
then threw unconditionally. Both were removed. The method now fails
fast if a non-football team is passed in, which matches the intent of a
football-specific season class.

\subsection{FootballTeam.addPlayer overload removed}

\code{FootballTeam} previously defined both
\code{addPlayer(Player)} (the overriding method) and
\code{addPlayer(\allowbreak FootballPlayer)} (a convenience overload that just cast
and re-dispatched). The overload was redundant because Java's overload
resolution picks the overriding method for any \code{FootballPlayer}
argument. It has been deleted.

\subsection{Demo renamed and decoupled}

The runnable demo is now \code{FootballConsoleDemo}. \code{Main} imports
it directly, so the \code{mvn exec:java} entry point no longer runs through
a class named ``Bootstrap''. The demo itself contains no football rules
and simply drives the application layer.

% ========================= 4. TESTING =========================
\section{Testing Strategy and Test Inventory}

The minimum requirement is twenty meaningful tests. The project ships
with \textbf{88}, spread across nineteen test classes. None of them are
null-only checks; each verifies a specific behaviour or invariant of
the code under test.

\subsection{Coverage by Package}

\begin{longtable}{lrl}
    \toprule
    \textbf{Test class} & \textbf{\#} & \textbf{Focus} \\
    \midrule
    \endhead
    \code{GameSessionTest}                  & 4  & Session wiring and progression state \\
    \code{FixtureTest}                      & 6  & Fixture construction and match attachment \\
    \code{LeagueTest}                       & 3  & Team and fixture membership invariants \\
    \code{FootballPlayerTest}               & 3  & Position, attribute profile, injury flow \\
    \code{FootballTeamTest}                 & 5  & Roster, lineup, tactic, training plan \\
    \code{FootballDataObjectsTest}          & 3  & Lineup, tactic, training plan invariants \\
    \code{FootballRulesetValidationTest}    & 7  & Lineup size, positions, duplicates, availability \\
    \code{FootballPointsCalculationTest}    & 2  & Points on win/draw/loss \\
    \code{FootballStandingsPolicyTest}      & 4  & Standings update, recomputation, ordering \\
    \code{RoundRobinFootballSchedulerTest}  & 4  & Fixture count, pairings, home/away balance \\
    \code{FootballSeasonTest}               & 6  & Weekly advance, injuries, lineup resolution, completion \\
    \code{FootballSportFactoryTest}         & 5  & League creation, team population, lineup generation, season wrap \\
    \code{FootballMatchTest}                & 6  & Match state, setup storage, replay rejection \\
    \code{FootballMatchEngineTest}          & 9  & Determinism, validation, stateless behaviour \\
    \code{FootballInjuryPolicyTest}         & 10 & Injury application, recovery, lineup rejection \\
    \code{EndToEndFootballFlowTest}         & 1  & Full season flow, bookkeeping invariants \\
    \code{InMemoryAssetProviderTest}        & 6  & Name, team, and logo data integrity \\
    \code{GameInitializationServiceTest}    & 3  & Session creation end-to-end through service \\
    \code{MainTest}                         & 1  & \code{Main.main} runs without throwing \\
    \midrule
    \textbf{Total}                          & \textbf{88} & \\
    \bottomrule
\end{longtable}

\subsection{What the Tests Actually Check}

A few examples of the kind of behaviour we verify:

\begin{itemize}
    \item \emph{Double round-robin invariants.} For a four-team league we
          assert exactly twelve fixtures, every team pair appears twice,
          no team plays itself, and each team has an equal number of home
          and away games.
    \item \emph{Standings bookkeeping.} The end-to-end test plays through
          a full season and asserts league-wide consistency: total goals
          for equals total goals against, total wins equals total losses,
          total draws is even, and the points total equals
          $3 \times \text{wins} + \text{draws}$. Policy-level tests also
          verify that \code{recordMatch} changes standings and does not
          double-count a match also present through a fixture.
    \item \emph{Match engine determinism.} Simulating the same pairing
          twice with a fresh engine yields the same score; simulating a
          different pairing in between does not shift the first result
          (the engine keeps no state).
    \item \emph{Replay prevention.} Calling \code{simulate} on a match
          that already has a result is a no-op; the stored result is
          untouched.
    \item \emph{Injury-availability link.} After a lost match, the losing
          team's stored lineup is cleared, one starter becomes unavailable,
          the ruleset rejects any lineup containing that player, and the
          recovery countdown is not decremented in the same played week.
          The player remains unavailable for the intended future match
          opportunities and becomes available at the correct later week.
    \item \emph{Football lineup policy.} Factory-generated and
          auto-resolved lineups are checked against the 1--4--4--2
          starting composition, and invalid positional compositions are
          rejected by the ruleset.
\end{itemize}

% ========================= 5. SIMPLIFICATIONS =========================
\section{Known Simplifications and Deferred Work}

A few pieces of the M1 design are modelled in code but are intentionally
not yet wired into the simulation loop.

\begin{itemize}
    \item \code{FootballAttributeProfile} (attack, defence, stamina,
          passing, speed) is attached to every generated player but the
          match engine does not yet consult it. Scores depend only on the
          deterministic seed.
    \item \code{FootballTrainingPlan} (focus, intensity, conditioning
          load, tactical load, recovery flag) is assigned to teams during
          generation, but the weekly ``Apply Training Plan'' step from
          Workflow~2 is not yet implemented. There are no training effects
          on player attributes or injury probability in M2.
    \item \code{FootballCoach} fields (specialisation, coaching rating)
          are stored but unused. Coaching bonuses will be introduced
          when training and attribute-driven simulation are added.
    \item The league format is fixed to a double round-robin. Cups, playoffs,
          and multi-division leagues are out of scope for M2.
    \item The ``future-season continuation hook'' mentioned in the task
          division is implemented at the stub level: \code{FootballSeason}
          exposes \code{canCreateNextSeason} and \code{createNextSeasonStub}
          which produce a fresh season carrying over the same teams. A full
          off-season flow is deferred.
    \item The user interface is not part of M2. Only the console demo
          is implemented.
\end{itemize}

% ========================= 6. RUN =========================
\section{How to Run the Project}

\subsection{Prerequisites}

\begin{itemize}
    \item Java 17 or newer
    \item Apache Maven 3.8 or newer.
\end{itemize}

\subsection{Commands}

From the project root:

\begin{lstlisting}[language=bash]
git clone https://github.com/mirzaisi/CE216_Project_Team_JavaIsSlow.git
cd CE216_Project_Team_JavaIsSlow

mvn compile
mvn test
mvn exec:java
\end{lstlisting}

The main class is declared in \code{pom.xml} via the exec plugin as
\code{com.playforgemanager.main.Main}.

\subsection{Expected Console Output}

Running \code{mvn exec:java} prints a short demo of a freshly initialised
football league. One week is simulated and the results, standings, and
next-week marker are printed. The exact output begins with:

\begin{lstlisting}
=== PlayForge Manager - M2 Demo ===
Sport: Football
League: PlayForge Demo League
Controlled Team: Red Hawks

Generated Teams
---------------
- Red Hawks       | Players: 19 | Coaches: 1
- Blue Wolves     | Players: 19 | Coaches: 1
- Golden Stars    | Players: 19 | Coaches: 1
- Iron Lions      | Players: 19 | Coaches: 1

Total fixtures generated: 12

Week 1 Results
----------------
...
\end{lstlisting}

Running \code{mvn test} should report
\code{Tests run: 88, Failures: 0, Errors: 0, Skipped: 0}.

% ========================= 7. RELEASE =========================
\section{Release Notes}

A GitHub Release is created from the final commit
that corresponds to this report. It contains:

\begin{itemize}
    \item The complete Maven project (source code, tests, \code{pom.xml}).
    \item This milestone report (PDF).
\end{itemize}

The repository URL, also the content of the submitted TXT file, is:

\begin{center}
    \url{https://github.com/mirzaisi/CE216_Project_Team_JavaIsSlow}
\end{center}

% ========================= 8. CONTRIBUTIONS =========================
\section{Team Contributions}

The division of work for M2 follows the updated task division for this
milestone. A short summary is given below for review convenience.

\begin{longtable}{ll}
    \toprule
    \textbf{Member}  & \textbf{M2 focus} \\
    \midrule
    \endhead
    Mirza İşi         & Project setup, shared framework, bootstrap flow, architecture check \\
    Sepehr Rahimi     & Football players, teams, ruleset, lineup and rules validation tests \\
    Doruk Faydalı     & Scheduler, standings, season progression, league integration pass   \\
    Çağan Parlapan    & Match engine, injury policy, asset provider, end-to-end verification, milestone report \\
    \bottomrule
\end{longtable}

% ========================= 9. AI DECLARATION =========================
\section*{Declaration on the Use of AI Assistance}
\addcontentsline{toc}{section}{Declaration on the Use of AI Assistance}

Consistent with the M1 declaration, AI-based tools were used in a
limited support role during M2. Specifically, they assisted with
structural review, wording of this report, and drafting of Java doc
comments. The design decisions, implementation, test content, and
submitted material remain the responsibility of the team.

\end{document}